\documentclass[11pt]{article}

\usepackage[margin=1in]{geometry}
\usepackage{amsmath,amssymb}
\usepackage[T1]{fontenc}

\title{Thresholded Sparse Covariance Estimation}
\author{}
\date{}

\begin{document}

\maketitle

Within each training fold, the two groups are centered separately using their
sample means. The centered observations are then pooled. If the pooled training
sample size is $n$, their cross-product matrix is divided by $n-2$, accounting
for the two estimated group means. The resulting pooled sample covariance is
denoted by $\widetilde\Sigma$.

We retain its diagonal entries. For $j\neq j'$, define
\[
\widehat\Sigma_{jj'}
=
\widetilde\Sigma_{jj'}
\mathbf 1
\left\{
|\widetilde\Sigma_{jj'}|
\geq
c\left[
\left(
\widetilde\Sigma_{jj}\widetilde\Sigma_{j'j'}
+\widetilde\Sigma_{jj'}^2
\right)
\frac{\log p}{n-2}
\right]^{1/2}
\right\}.
\]
The factor inside the threshold accounts for entry-specific scale under the
Gaussian simulation model. The diagnostic comparison uses
$c\in\{1,1.25,1.5\}$; we use $c=1.25$ as the default.

If the thresholded matrix has smallest eigenvalue below a fixed positive
numerical tolerance, we add the minimum diagonal shift required to reach that
tolerance. This operation preserves the eigenvectors and all eigenvalue gaps.

The implementation is contained in
\texttt{src/002\_threshold\_sparse\_covariance.R}. The validation script
\texttt{code/002\_validate\_sparse\_covariance\_estimator.R} reads the two
population covariance matrices and saves repetition-level operator-norm and
leading-eigenvector diagnostics in
\texttt{data/intermediate/threshold\_covariance\_operator\_norm\_diagnostics.rds}.

\end{document}
